\chapter{Układy z wymuszeniem i funkcje przejścia}
\label{ch:lin28-transfer}
\index{funkcja przejścia}
\index{odpowiedź impulsowa}
\index{układ liniowy!wymuszenie}
\index{realizacja!funkcji wymiernej}

\noindent\textbf{Cel rozdziału.}
W rozdziałach~\ref{ch:lin25-control} i~\ref{ch:lin27-realizations}
opisaliśmy sterowalność, obserwowalność oraz minimalne realizacje
za pomocą parametrów Markowa. Zbadamy teraz, jak macierze
stanu wyznaczają odpowiedź na \emph{dowolne wymuszenie},
oraz jak tę odpowiedź koduje \emph{macierzowa funkcja przejścia}.
Pokażemy obie wersje czasowe, skonstruujemy realizację
zadanej skalarnej funkcji wymiernej i wyjaśnimy, dlaczego
stabilność odpowiedzi wejście--wyjście nie zawsze pociąga
za sobą stabilność wszystkich zmiennych stanu.

W całym rozdziale przyjmujemy
\(n,m,p\ge1\), \(A\in\mathbb C^{n\times n}\),
\(B\in\mathbb C^{n\times m}\),
\(C\in\mathbb C^{p\times n}\) i
\(D\in\mathbb C^{p\times m}\).
W normach wektorowych stosujemy normę euklidesową,
a w macierzowych odpowiadającą jej normę operatorową.
Symbol \(I\) oznacza macierz jednostkową odpowiedniego rozmiaru.

\section{Rozwiązanie z warunkiem początkowym}
\label{sec:lin28-variation}

\begin{twierdzenie}[Wzór wariacji stałych]
\label{thm:lin28-variation}
Niech \(T>0\), \(x_0\in\mathbb C^n\), a
\(u\colon[0,T]\to\mathbb C^m\) będzie funkcją ciągłą.
Zagadnienie
\[
 x'(t)=Ax(t)+Bu(t),\qquad x(0)=x_0
\]
ma dokładnie jedno rozwiązanie klasy \(C^1\).
Dla każdego \(t\in[0,T]\) jest nim
\[
 x(t)=e^{tA}x_0+
        \int_0^t e^{(t-\tau)A}Bu(\tau)\,\mathrm d\tau.
 \tag{28.1}\label{eq:lin28-variation}
\]
Jeżeli \(y(t)=Cx(t)+Du(t)\), to
\[
 y(t)=Ce^{tA}x_0+Du(t)+
       \int_0^t Ce^{(t-\tau)A}Bu(\tau)\,\mathrm d\tau.
 \tag{28.2}\label{eq:lin28-output}
\]
\end{twierdzenie}
\begin{proof}
Dla dowolnego rozwiązania połóżmy
\(v(t)=e^{-tA}x(t)\).
Macierze \(A\) i \(e^{-tA}\) komutują, ponieważ
wykładnik jest granicą wielomianów w \(A\).
Reguła różniczkowania iloczynu daje
\[
 v'(t)=-Ae^{-tA}x(t)+e^{-tA}(Ax(t)+Bu(t))
      =e^{-tA}Bu(t).
\]
Po scałkowaniu od \(0\) do \(t\), skorzystaniu
z \(v(0)=x_0\) i pomnożeniu przez \(e^{tA}\)
otrzymujemy wzór~\eqref{eq:lin28-variation};
zastosowaliśmy przy tym
\(e^{tA}e^{-\tau A}=e^{(t-\tau)A}\).
Zatem każde rozwiązanie musi mieć wskazaną postać.

Odwrotnie, całkowana funkcja jest ciągła
w obu zmiennych, więc prawą stronę
\eqref{eq:lin28-variation} można różniczkować
względem \(t\) na mocy reguły Leibniza.
Pochodna całki jest równa
\[
 Bu(t)+\int_0^t Ae^{(t-\tau)A}Bu(\tau)\,\mathrm d\tau.
\]
Po dodaniu pochodnej \(e^{tA}x_0\)
dostajemy \(x'=Ax+Bu\);
dla \(t=0\) całka znika, zatem \(x(0)=x_0\).
Istnienie i jednoznaczność zostały udowodnione.
Wstawienie wzoru na \(x\) do \(y=Cx+Du\)
daje~\eqref{eq:lin28-output}.
\end{proof}

\begin{uwaga}[Wymuszenia całkowalne]
\label{rem:lin28-integrable}
Jeśli \(u\in L^1([0,T];\mathbb C^m)\), całka
w~\eqref{eq:lin28-variation} nadal istnieje.
Wzór definiuje jedyne rozwiązanie absolutnie ciągłe,
spełniające równanie stanu prawie wszędzie
oraz warunek początkowy.
Istotnie, funkcja \(v(t)=x_0+\int_0^t e^{-\tau A}Bu(\tau)
\,\mathrm d\tau\) jest absolutnie ciągła;
stosuje się do niej regułę iloczynu prawie wszędzie,
a dowód jednoznaczności przebiega jak wyżej.
Jest to użyteczne dla wymuszeń przedziałami stałych.
\end{uwaga}

\begin{twierdzenie}[Rozwiązanie rekurencji z wymuszeniem]
\label{thm:lin28-discrete}
Dla dowolnego \(x_0\in\mathbb C^n\)
oraz dowolnego ciągu \(u_k\in\mathbb C^m\),
\(k\ge0\), układ
\(x_{k+1}=Ax_k+Bu_k\), \(y_k=Cx_k+Du_k\)
ma dokładnie jedno rozwiązanie. Dla \(k\ge0\)
zachodzi
\[
 \begin{aligned}
 x_k&=A^kx_0+\sum_{j=0}^{k-1}A^{k-1-j}Bu_j,\\
 y_k&=CA^kx_0+Du_k+
       \sum_{j=0}^{k-1}CA^{k-1-j}Bu_j.
 \end{aligned}
 \tag{28.3}\label{eq:lin28-discrete}
\]
Suma z górnym indeksem \(k-1=-1\)
jest z definicji pusta i równa zeru.
\end{twierdzenie}
\begin{proof}
Dla \(k=0\) obie równości są natychmiastowe.
Jeśli pierwsza zachodzi dla \(k\), to
\[
 x_{k+1}
 =A^{k+1}x_0+
   \sum_{j=0}^{k-1}A^{k-j}Bu_j+Bu_k
 =A^{k+1}x_0+\sum_{j=0}^{k}A^{k-j}Bu_j.
\]
Otrzymujemy indukcyjnie pierwszą równość
dla wszystkich \(k\). Druga wynika przez
pomnożenie przez \(C\) i dodanie \(Du_k\).
Rekurencja jednoznacznie wyznacza każdy kolejny stan.
\end{proof}

\section{Odpowiedź impulsowa i splot}
\label{sec:lin28-impulse}

\begin{definicja}[Jądro odpowiedzi i człon bezpośredni]
\label{def:lin28-kernel}
Dla układu ciągłego definiujemy
\(h(t)=Ce^{tA}B\) dla \(t\ge0\).
W układzie dyskretnym określamy
\(h_0=D\) oraz \(h_k=CA^{k-1}B\) dla \(k\ge1\).
Macierz \(D\) nazywamy
\emph{członem bezpośrednim} (przeniesieniem bezpośrednim).
\end{definicja}

\begin{twierdzenie}[Odpowiedź przy zerowym stanie początkowym]
\label{thm:lin28-convolution}
Jeżeli \(x_0=0\), to dla ciągłego, ciągle wymuszanego
układu i dla dyskretnego układu zachodzą odpowiednio
\[
 y(t)=Du(t)+\int_0^t h(t-\tau)u(\tau)\,\mathrm d\tau,
 \qquad
 y_k=\sum_{j=0}^{k}h_{k-j}u_j.
 \tag{28.4}\label{eq:lin28-convolution}
\]
W szczególności przy wejściu dyskretnym
\(u_0=w\in\mathbb C^m\), \(u_k=0\) dla \(k\ge1\),
odpowiedź ma postać
\(y_0=Dw\), \(y_k=CA^{k-1}Bw\) dla \(k\ge1\).
\end{twierdzenie}
\begin{proof}
Pierwsza formuła to
\eqref{eq:lin28-output} przy \(x_0=0\)
po podstawieniu \(h(t)=Ce^{tA}B\).
Z~\eqref{eq:lin28-discrete} dostajemy
\(y_k=Du_k+\sum_{j=0}^{k-1}h_{k-j}u_j\),
co jest drugą formułą, ponieważ \(h_0=D\).
Dla wskazanego wejścia impulsowego
w sumie pozostaje tylko składnik z \(j=0\).
\end{proof}

\begin{uwaga}[Impuls Diraca i skok wyjścia]
\label{rem:lin28-dirac}
Formalny zapis ciągłej odpowiedzi impulsowej
to \(D\delta_0+h\), gdzie \(\delta_0\)
jest deltą Diraca. Nie jest to zwykła funkcja,
jeśli \(D\ne0\). Dokładne znaczenie podają
wzory~\eqref{eq:lin28-convolution}:
człon \(Du(t)\) działa natychmiastowo,
a pozostała część jest całką splotową.
Nie musimy definiować dystrybucji,
aby posługiwać się tymi wzorami.
\end{uwaga}

\begin{przyklad}[Wymuszenie jednostkowe i stan początkowy]
\label{prz:lin28-step}
Niech \(n=m=p=1\), \(A=-2\), \(B=3\), \(C=1\),
\(D=0\), \(x_0=4\), a \(u(t)=1\) dla \(t\ge0\).
Wzór~\eqref{eq:lin28-output} daje
\[
 y(t)=4e^{-2t}+3\int_0^t e^{-2(t-\tau)}
 \,\mathrm d\tau
 =4e^{-2t}+\frac32(1-e^{-2t}).
\]
Całkę obliczyliśmy podstawiając \(s=t-\tau\):
\(\int_0^t e^{-2s}\,\mathrm ds=(1-e^{-2t})/2\).
Granica \(y(t)\to3/2\) wynika z \(e^{-2t}\to0\).
Widać oddzielnie zanikającą odpowiedź od stanu
początkowego i odpowiedź na wymuszenie.
\end{przyklad}

\section{Transmitancja ciągła i dyskretna}
\label{sec:lin28-transfer}

\begin{definicja}[Macierzowa funkcja przejścia]
\label{def:lin28-transfer}
Dla opisanej na początku rozdziału
czwórki macierzy \((A,B,C,D)\)
definiujemy funkcję wymierną
\[
 G(\zeta)=D+C(\zeta I-A)^{-1}B,
 \qquad \zeta\in\mathbb C\setminus\sigma(A).
 \tag{28.5}\label{eq:lin28-transfer}
\]
W kontekście ciągłym argument oznaczamy przez \(s\),
a w dyskretnym przez \(z\).
Wartości funkcji leżą w \(\mathbb C^{p\times m}\).
W punktach \(\sigma(A)\) poszczególne elementy
macierzy \(G\) mogą mieć przedłużenia holomorficzne;
nie zakładamy, że każdy taki punkt jest biegunem \(G\).
\end{definicja}

\begin{twierdzenie}[Transformata Laplace'a odpowiedzi]
\label{thm:lin28-laplace}
Niech \(u\colon[0,\infty)\to\mathbb C^m\)
będzie funkcją przedziałami ciągłą o zwartym nośniku,
a \(x_0=0\). Dla każdego \(s\in\mathbb C\)
o dostatecznie dużej części rzeczywistej
istnieją całki
\[
 U(s)=\int_0^\infty e^{-st}u(t)\,\mathrm dt,
 \qquad
 Y(s)=\int_0^\infty e^{-st}y(t)\,\mathrm dt,
\]
oraz zachodzi równość
\[
 Y(s)=G(s)U(s).
 \tag{28.6}\label{eq:lin28-laplace}
\]
\end{twierdzenie}
\begin{proof}
Ponieważ macierz \(A\) jest skończona,
z szeregu wykładniczego wynika oszacowanie
\(\|e^{tA}\|\le e^{t\|A\|}\) dla \(t\ge0\).
Wzór~\eqref{eq:lin28-variation} i zwartość nośnika
\(u\) dają \(\|x(t)\|\le K e^{t\|A\|}\)
dla pewnego \(K>0\). Zatem przy
\(\operatorname{Re}s>\|A\|\) możemy zastosować
transformatę Laplace'a i całkowanie przez części:
\[
 \int_0^\infty e^{-st}x'(t)\,\mathrm dt
   =sX(s)-x(0)=sX(s).
\]
Wyraz brzegowy w nieskończoności zanika
na mocy podanego oszacowania.
Z równania stanu dostajemy
\((sI-A)X(s)=BU(s)\).
Macierz \(sI-A\) jest odwracalna,
gdy \(\operatorname{Re}s>\|A\|\), gdyż
wartości własne \(A\) mają moduł
nie większy od \(\|A\|\).
Stąd \(X(s)=(sI-A)^{-1}BU(s)\).
Transformacja \(y=Cx+Du\) prowadzi do
\(Y(s)=(C(sI-A)^{-1}B+D)U(s)\),
co kończy dowód.
\end{proof}

\begin{twierdzenie}[Transformata \(Z\) odpowiedzi]
\label{thm:lin28-z}
Niech \(u_k\in\mathbb C^m\) ma nośnik skończony,
a \(x_0=0\). Wówczas dla wszystkich zespolonych
\(z\) o dostatecznie dużym module
szeregi
\[
 U(z)=\sum_{k=0}^\infty u_kz^{-k},
 \qquad
 Y(z)=\sum_{k=0}^\infty y_kz^{-k}
\]
są zbieżne i zachodzi \(Y(z)=G(z)U(z)\).
Ponadto dla \(|z|>\|A\|\)
\[
 G(z)=D+\sum_{j=0}^\infty CA^jB\,z^{-j-1}.
 \tag{28.7}\label{eq:lin28-laurent}
\]
\end{twierdzenie}
\begin{proof}
Przy skończenie wielu niezerowych wejściach
\eqref{eq:lin28-discrete} daje
\(\|x_k\|\le K(1+\|A\|)^k\) z pewnym \(K>0\);
stąd szeregi dla \(X(z)=\sum x_k z^{-k}\)
i \(Y(z)\) są zbieżne przy dostatecznie dużym \(|z|\).
Ponieważ \(x_0=0\), bezwzględna zbieżność
umożliwia przesunięcie indeksu:
\[
 \sum_{k\ge0}x_{k+1}z^{-k}
 =z\sum_{k\ge0}x_kz^{-k}=zX(z).
\]
Po transformacji rekurencji otrzymujemy
\((zI-A)X(z)=BU(z)\) i
\(Y(z)=CX(z)+DU(z)=G(z)U(z)\).

Dla \(|z|>\|A\|\) szereg geometryczny macierzy
\(\sum_{j\ge0}(A/z)^j\)
jest bezwzględnie zbieżny i stanowi odwrotność
\(I-A/z\): po pomnożeniu sumy częściowej
przez \(I-A/z\) reszta wynosi \((A/z)^{N+1}\)
i dąży do zera. Dlatego
\((zI-A)^{-1}=\sum_{j\ge0} A^j z^{-j-1}\).
Podstawienie do definicji \(G\)
prowadzi do~\eqref{eq:lin28-laurent}.
\end{proof}

\begin{wniosek}[Parametry Markowa z transmitancji]
\label{cor:lin28-markov}
Dla \(i=1,2\) niech \(n_i\ge1\),
\(A_i\in\mathbb C^{n_i\times n_i}\),
\(B_i\in\mathbb C^{n_i\times m}\),
\(C_i\in\mathbb C^{p\times n_i}\),
\(D_i\in\mathbb C^{p\times m}\), gdzie \(m,p\ge1\).
Ich wymierne funkcje przejścia są jednakowe
dokładnie wtedy, gdy \(D_1=D_2\) oraz
\(C_1A_1^kB_1=C_2A_2^kB_2\)
dla każdego całkowitego \(k\ge0\).
Równoważnie ich odpowiedzi na wszystkie
wejścia przy zerowym stanie początkowym
są jednakowe, zarówno w modelu ciągłym,
jak i dyskretnym.
\end{wniosek}
\begin{proof}
Z równości funkcji wymiernych wynika równość
ich rozwinięć Laurenta przy nieskończoności.
Przez~\eqref{eq:lin28-laurent} mamy wtedy
równość członów bezpośrednich i wszystkich
parametrów Markowa. Odwrotnie, równość
współczynników oznacza, że funkcje \(G_i\)
są identyczne na wspólnym obszarze \(|z|>R\);
jako funkcje wymierne są więc identyczne
wszędzie, gdzie obie są określone.
Równoważność z odpowiedziami czasowymi
wynika z~\eqref{eq:lin28-convolution}
i szeregu \(Ce^{tA}B=\sum_{k\ge0}t^kCA^kB/k!\).
Dla odwrotnej implikacji wystarczy
użyć dyskretnego wejścia impulsowego:
jego kolejne wyjścia odczytują \(D\)
oraz wszystkie \(CA^kB\).
\end{proof}

\section{Bieguny, wartości własne i realizacje}
\label{sec:lin28-poles}

\begin{stwierdzenie}[Wymierność i możliwe bieguny]
\label{prop:lin28-rational}
Każdy współczynnik macierzy \(G\)
jest funkcją wymierną właściwą:
jego granica przy \(|\zeta|\to\infty\)
jest równa odpowiedniemu współczynnikowi \(D\).
Każdy biegun elementu \(G\) jest
wartością własną \(A\), ale odwrotna
implikacja bez dodatkowych założeń jest fałszywa.
\end{stwierdzenie}
\begin{proof}
Dla \(\det(\zeta I-A)\ne0\) stosujemy
wzór na macierz odwrotną:
\((\zeta I-A)^{-1}
 =\operatorname{adj}(\zeta I-A)/
  \det(\zeta I-A)\).
Mianownik jest wielomianem monicznym
stopnia \(n\), a elementy macierzy dołączonej
są wielomianami stopnia co najwyżej \(n-1\).
Po pomnożeniu przez stałe macierze \(B,C\)
otrzymujemy funkcję ściśle właściwą;
dodanie \(D\) daje funkcję właściwą
z granicą \(D\). Bieguny mogą więc leżeć
tylko wśród zer \(\det(\zeta I-A)\),
czyli wartości własnych \(A\).
Możliwość skrócenia wspólnego czynnika
licznika i mianownika wykaże
Przykład~\ref{prz:lin28-hidden-mode}.
\end{proof}

\begin{przyklad}[Niewidoczny niestabilny stan]
\label{prz:lin28-hidden-mode}
Niech
\[
 A=\begin{pmatrix}1&0\\0&-1\end{pmatrix},
 \quad B=\begin{pmatrix}0\\1\end{pmatrix},
 \quad C=\begin{pmatrix}0&1\end{pmatrix},\quad D=0.
\]
Dla \(\zeta\notin\{1,-1\}\)
mamy \((\zeta I-A)^{-1}
=\operatorname{diag}((\zeta-1)^{-1},(\zeta+1)^{-1})\).
Stąd \(G(\zeta)=1/(\zeta+1)\):
wartość własna \(1\) nie jest biegunem
funkcji przejścia. Pierwszego stanu
ani nie można osiągnąć z wejścia,
ani nie widać go na wyjściu.
Mimo że odpowiedź na każde ograniczone
wymuszenie jest ograniczona,
niezerowy początkowy pierwszy stan
rośnie jak \(e^t\).
Realizacja jednowymiarowa
\(A_{\min}=-1\), \(B_{\min}=C_{\min}=1\), \(D=0\)
daje tę samą funkcję przejścia.
\end{przyklad}

\begin{twierdzenie}[Realizacja skalarnej funkcji właściwej]
\label{thm:lin28-companion}
Niech \(n\ge1\) i niech
\[
 p(\zeta)=\zeta^n+a_{n-1}\zeta^{n-1}
      +\cdots+a_1\zeta+a_0,\qquad
 q(\zeta)=b_{n-1}\zeta^{n-1}
      +\cdots+b_1\zeta+b_0
\]
będą wielomianami o współczynnikach zespolonych,
przy czym \(p\) jest moniczny i \(\deg q<n\)
(dopuszczamy \(q=0\)).
Dla dowolnego \(d\in\mathbb C\)
funkcja \(g(\zeta)=d+q(\zeta)/p(\zeta)\)
posiada realizację wymiaru \(n\):
\[
 A=\begin{pmatrix}
 0&1&0&\cdots&0\\
 0&0&1&\ddots&\vdots\\
 \vdots&&\ddots&\ddots&0\\
 0&\cdots&0&0&1\\
 -a_0&-a_1&\cdots&-a_{n-2}&-a_{n-1}
 \end{pmatrix},
 \quad B=e_n,\quad
 C=\begin{pmatrix}b_0&b_1&\cdots&b_{n-1}\end{pmatrix},
 \quad D=d.
 \tag{28.8}\label{eq:lin28-companion}
\]
Przy \(n=1\) należy rozumieć \(A=(-a_0)\).
Para \((A,B)\) jest zawsze sterowalna.
Realizacja jest minimalna dokładnie wtedy,
gdy \(p\) i \(q\) nie mają wspólnego pierwiastka
w \(\mathbb C\), czyli są względnie pierwsze.
\end{twierdzenie}
\begin{proof}
Oznaczmy \(v(\zeta)=(1,\zeta,\ldots,\zeta^{n-1})^T\).
Bezpośrednie mnożenie macierzy w
\eqref{eq:lin28-companion} daje
\[
 (\zeta I-A)v(\zeta)=p(\zeta)e_n:
\]
pierwsze \(n-1\) współrzędnych są zerowe,
a ostatnia to
\(a_0+a_1\zeta+\cdots+a_{n-1}\zeta^{n-1}+\zeta^n\).
Jeżeli \(p(\zeta)\ne0\), to
\((\zeta I-A)^{-1}B=v(\zeta)/p(\zeta)\),
a zatem \(C(\zeta I-A)^{-1}B=q(\zeta)/p(\zeta)\).
Daje to żądaną funkcję \(g\).

Dalej, \(Ae_j=e_{j-1}-a_{j-1}e_n\)
dla \(j=2,\ldots,n\). Stąd
\(B,AB,\ldots,A^{n-1}B\),
zapisane w odwrotnie uporządkowanej bazie
\(e_n,e_{n-1},\ldots,e_1\),
tworzą macierz trójkątną z jedynkami
na przekątnej: w \(A^kB\) współczynnik
przy \(e_{n-k}\) jest równy \(1\),
a przy \(e_j\) dla \(j<n-k\) jest zerowy.
Macierz sterowalności ma więc rangę \(n\).

W celu sprawdzenia obserwowalności
stosujemy kryterium PBH z
Wniosku~\ref{cor:lin25-pbh-observe}.
Każda wartość własna \(\lambda\) jest zerem
\(p\), a wektor \(v(\lambda)\ne0\)
spełnia \(Av(\lambda)=\lambda v(\lambda)\).
Ponadto każde rozwiązanie
\((A-\lambda I)w=0\) jest wielokrotnością
\(v(\lambda)\): pierwsze \(n-1\) wierszy
wymuszają \(w_{j+1}=\lambda w_j\).
Warunek obserwowalności jest więc równoważny
\(Cv(\lambda)=q(\lambda)\ne0\)
dla każdego pierwiastka \(p\).
Nad \(\mathbb C\) jest to równoważne
względnej pierwszości \(p,q\).
Ponieważ sterowalność już udowodniono,
Twierdzenie~\ref{thm:lin27-hankel-rank}
o minimalnym wymiarze realizacji
kończy dowód.
\end{proof}

\begin{uwaga}[Funkcja stała i stopień minimalny]
\label{rem:lin28-constant}
Jeśli \(q=0\), wzór nadal daje realizację,
lecz nie jest ona minimalna: funkcję
stałą \(g=d\) realizuje układ bez stanów,
z samym członem bezpośrednim \(D=d\).
Jeżeli \(q/p\) skraca się do
\(\widetilde q/\widetilde p\) z
\(\gcd(\widetilde q,\widetilde p)=1\),
to powyższa konstrukcja dla
\(\widetilde p,\widetilde q\)
daje minimalny wymiar
\(\deg\widetilde p\); w stopniu \(0\)
otrzymujemy realizację bez stanów.
Wykazaliśmy to dla funkcji skalarnych;
dla funkcji macierzowych samo zliczanie
biegunów poszczególnych elementów nie
zastępuje rangi macierzy Hankela
z Twierdzenia~\ref{thm:lin27-hankel-rank}.
\end{uwaga}

\begin{przyklad}[Realizacja i odpowiedź impulsowa]
\label{prz:lin28-companion}
Niech \(g(s)=(2s+3)/(s^2+3s+2)\).
Konstrukcja z
Twierdzenia~\ref{thm:lin28-companion} daje
\[
 A=\begin{pmatrix}0&1\\-2&-3\end{pmatrix},
 \qquad B=\begin{pmatrix}0\\1\end{pmatrix},
 \qquad C=\begin{pmatrix}3&2\end{pmatrix},\qquad D=0.
\]
Istotnie, \((sI-A)^{-1}B
=(1,s)^T/(s^2+3s+2)\).
Ponieważ \(p(s)=(s+1)(s+2)\),
a \(q(-1)=1\), \(q(-2)=-1\),
realizacja jest minimalna.
Rozkład na ułamki proste ma postać
\[
 \frac{2s+3}{(s+1)(s+2)}
 =\frac1{s+1}+\frac1{s+2}.
\]
Jądro odpowiedzi impulsowej jest zatem
\(h(t)=e^{-t}+e^{-2t}\) dla \(t\ge0\):
potwierdza to jednoznaczność transformaty
Laplace'a, można też bezpośrednio
obliczyć \(Ce^{tA}B\).
Dla wejścia \(u(t)=1\) i \(x_0=0\) otrzymujemy
\[
 y(t)=\int_0^t(e^{-(t-\tau)}+e^{-2(t-\tau)})
 \,\mathrm d\tau
 =\frac32-e^{-t}-\frac12e^{-2t}.
\]
W szczególności \(y(0)=0\)
i \(y(t)\to3/2=g(0)\), zgodnie
z bezpośrednim obliczeniem całki.
\end{przyklad}

\section{Oszacowania wejście--wyjście}
\label{sec:lin28-bounds}

\begin{twierdzenie}[Ograniczone wejście i stabilny układ ciągły]
\label{thm:lin28-bibo}
Niech każda wartość własna
\(A\in\mathbb C^{n\times n}\) ma ujemną
część rzeczywistą.
Istnieją stałe \(M,\alpha>0\)
takie, że dla wszystkich \(t\ge0\)
\(\|e^{tA}\|\le Me^{-\alpha t}\).
Dla dowolnego \(x_0\in\mathbb C^n\)
i ograniczonej funkcji
\(u\in L^\infty([0,\infty);\mathbb C^m)\)
odpowiedź wyznaczona wzorem
\eqref{eq:lin28-output} spełnia
dla prawie wszystkich \(t\ge0\)
\[
 \|y(t)\|\le
 M\|C\|e^{-\alpha t}\|x_0\|
 +\left(\|D\|+\frac{M\|C\|\|B\|}{\alpha}\right)
     \|u\|_{L^\infty}.
 \tag{28.9}\label{eq:lin28-bibo}
\]
\end{twierdzenie}
\begin{proof}
Rozkład Jordana (Rozdział~\ref{ch:lin17-jordan})
wyraża \(e^{tA}\), po zmianie bazy,
jako sumę bloków zawierających
\(e^{\lambda t}t^j/j!\)
dla \(0\le j<n\). Jeżeli
\(\beta=-\max_{\lambda\in\sigma(A)}
\operatorname{Re}\lambda>0\),
to dla \(\alpha=\beta/2\)
wszystkie funkcje
\(t^je^{-\beta t/2}\) są ograniczone
na \([0,\infty)\).
Po uwzględnieniu norm macierzy zmiany
bazy dostajemy wspólne \(M>0\)
takie, że \(\|e^{tA}\|\le Me^{-\alpha t}\).

Dla całkowej części~\eqref{eq:lin28-output}
nierówność trójkąta i podmultiplikatywność
normy dają
\[
 \left\|\int_0^t Ce^{(t-\tau)A}Bu(\tau)
 \,\mathrm d\tau\right\|
 \le M\|C\|\|B\|\|u\|_{L^\infty}
       \int_0^t e^{-\alpha(t-\tau)}\,\mathrm d\tau.
\]
Ostatnia całka jest równa
\((1-e^{-\alpha t})/\alpha\le1/\alpha\).
Pozostałe składniki odpowiedzi szacujemy
przez \(M\|C\|e^{-\alpha t}\|x_0\|\)
oraz \(\|D\|\|u\|_{L^\infty}\).
Dodanie tych oszacowań dowodzi tezy.
\end{proof}

\begin{wniosek}[Dyskretne oszacowanie]
\label{cor:lin28-discrete-bibo}
Jeżeli \(\rho(A)<1\), to istnieją
\(M>0\), \(r\in(0,1)\) takie, że
\(\|A^k\|\le Mr^k\) dla wszystkich
\(k\ge0\). Dla każdego ograniczonego
ciągu wejść i każdego stanu początkowego
zachodzi
\[
 \|y_k\|\le
 M\|C\|r^k\|x_0\|+
 \left(\|D\|+\frac{M\|C\|\|B\|}{1-r}\right)
 \sup_{j\ge0}\|u_j\|.
 \tag{28.10}\label{eq:lin28-discrete-bibo}
\]
\end{wniosek}
\begin{proof}
Dla każdego \(r\) z \(\rho(A)<r<1\)
rozkład Jordana i ograniczoność
\(k^j(\rho(A)/r)^k\) (dla niezerowych
wartości własnych) dają oszacowanie
\(\|A^k\|\le Mr^k\); bloki nilpotentne
znikają dla dostatecznie dużych \(k\).
Wstawiamy je do~\eqref{eq:lin28-discrete}
i korzystamy z
\(\sum_{j=0}^{k-1}r^{k-1-j}
  =(1-r^k)/(1-r)\le1/(1-r)\).
\end{proof}

\begin{uwaga}[Zakres wniosków o stabilności]
\label{rem:lin28-hidden-stability}
Twierdzenie~\ref{thm:lin28-bibo} oraz
Wniosek~\ref{cor:lin28-discrete-bibo}
są kryteriami \emph{dostatecznymi}.
Bez sterowalności i obserwowalności
nie wolno wywodzić stabilności całej
macierzy \(A\) ze stabilności transmitancji:
Przykład~\ref{prz:lin28-hidden-mode}
pokazuje przeciwną sytuację.
W realizacji minimalnej ukryty tryb
nie może wystąpić, co tłumaczy rolę
kryteriów Kalmana i macierzy Hankela.
\end{uwaga}

\section{Podsumowanie i powiązania}
\label{sec:lin28-summary}
Dynamikę z warunkiem początkowym
i wymuszeniem opisują dokładne wzory
\eqref{eq:lin28-output} i
\eqref{eq:lin28-discrete}. Dla zerowego
stanu początkowego odpowiedzi są splotami,
a ich transformaty są mnożeniem przez
wymierną macierz \(G\).
Parametry Markowa stanowią współczynniki
rozwinięcia \(G\) przy nieskończoności;
minimalny wymiar realizacji jest związany
z rangą macierzy Hankela z
Rozdziału~\ref{ch:lin27-realizations}.
W kolejnym rozdziale wrócimy
do geometrii podprzestrzeni,
zaczynając od grassmannianów.
